\documentclass[11pt]{article}
\usepackage{amsmath, amsthm, amssymb}
\usepackage{hyperref}
\usepackage{geometry}
\usepackage{setspace}

\geometry{margin=1in}

\newtheorem{theorem}{Theorem}
\newtheorem{lemma}{Lemma}
\newtheorem{corollary}{Corollary}

\newcommand{\R}{\mathbb{R}}
\newcommand{\T}{\mathcal{T}}
\newcommand{\W}{\mathcal{W}}
\newcommand{\lhat}{\hat{\lambda}}
\newcommand{\ghat}{\hat{\gamma}}
\newcommand{\muhat}{\hat{\mu}}
\newcommand{\phihat}{\hat{\varphi}}
\newcommand{\Op}{\mathcal{O}_P}
\newcommand{\id}{\mathrm{id}}
\newcommand{\sups}{\sup_{t \in \T}}
\newcommand{\norm}[1]{\|#1\|_\infty}
\newcommand{\TV}{\mathrm{TV}}

\title{Asymptotic Consistency of a Fr\'echet Mean Estimator for the Latent Deformation Model}
\author{Cody Carroll}
\date{Nov. 2026}

\begin{document}
\maketitle
\onehalfspacing

\section{Introduction}

Carroll and M\"{u}ller (2023) propose the latent deformation model (LDM)
\begin{equation}
  X_{ij}(t) = A_{ij}\,(\lambda \circ \Psi_j \circ H_i)(t),
  \qquad i = 1,\ldots,n,\quad j = 1,\ldots,p,
  \label{eq:LDM}
\end{equation}
where $\lambda$ is the latent curve with $\norm{\lambda} = 1$, $A_{ij} > 0$ are random
amplitude factors, $\Psi_j$ are fixed component-level deformations, and $H_i$ are random
subject-level deformations. Without loss of generality we take $\T = [0,1]$. The
deformations satisfy the standardization conditions
\begin{align}
  \frac{1}{p}\sum_{j=1}^p \Psi_j^{-1}(t) &= t \quad \text{for all } t \in \T,
  \label{eq:std-psi}\\
  E\, H_i^{-1}(t) &= t \quad \text{for all } t \in \T.
  \label{eq:std-H}
\end{align}
The $j$th component tempo is $\gamma_j = \lambda \circ \Psi_j$, so that the normalized curves are
$X_{ij}/\norm{X_{ij}} = \gamma_j \circ H_i$.

Carroll and M\"{u}ller (2023) estimate the component tempos as follows. For each component
$j$, the normalized curves $\{X_{ij}/\norm{X_{ij}}\}_{i=1}^n$ are aligned by pairwise warping
(Tang and M\"{u}ller, 2008) with penalty parameter $\eta_1$, which yields estimates
$\tilde{H}_i^{(j)}$ of the subject-level deformations. These are averaged over components and
used to align each component:
\begin{equation}
  \hat{H}_i = \frac{1}{p}\sum_{j=1}^p \tilde{H}_i^{(j)},
  \qquad
  \ghat_j = \frac{1}{n}\sum_{i=1}^n \frac{X_{ij}\circ\hat{H}_i^{-1}}{\norm{X_{ij}}},
  \qquad j = 1,\ldots,p.
  \label{eq:gammahat}
\end{equation}
When the curves are observed with noise on a grid of $m$ time points, kernel-smoothed curves
with bandwidth $b$ are substituted for $X_{ij}$. The latent curve is estimated by a global
alignment estimator: one component curve is selected at random for each subject, these $n$
normalized curves are aligned by pairwise warping, and the aligned curves are averaged.

The convergence results of Carroll and M\"{u}ller (2023) are obtained under the following
assumptions on (L) the components of the LDM and (S) the smoothing step.
\begin{itemize}
  \item[(L1)] The latent curve $\lambda \in C^2(\T)$ is a bounded function. For any
    nondegenerate interval $\T_0 \subset \T$, $0 < \int_{\T_0}\lambda'(t)^2\,dt < \infty$.
  \item[(L2)] For $j = 1,\ldots,p$, $\sup_{1\le i\le n} A_{ij} = \Op(1)$ and
    $\sup_{1\le i\le n} A_{ij}^{-1} = \Op(1)$.
  \item[(S0)] The time points $t_1,\ldots,t_m$ depend on the sample size, $m = m(n)$, and
    constitute a dense regular design generated by a smooth design density $f$ with
    $\inf_{t\in\T} f(t) > 0$ and bounded second derivative, according to
    $t_s = F^{-1}\bigl(\tfrac{s-1}{m-1}\bigr)$, $s = 1,\ldots,m$, where $F^{-1}$ is the
    quantile function associated with $f$.
  \item[(S1)] The kernel function $K$ is a probability density function with support
    $[-1,1]$, symmetric around zero, and uniformly continuous on its support, with
    $\int_{-1}^{1}K^2(u)\,du < \infty$.
  \item[(S2)] The sequences $m = m(n)$ and $b = b(n)$ satisfy $0 < b < \infty$, and
    $m \to \infty$, $b \to 0$, and $mb^2(\log b)^{-1} \to \infty$ as $n \to \infty$.
\end{itemize}
Under (L1), (L2), and (S0--S2), Theorem 1 of Carroll and M\"{u}ller (2023) gives
\begin{equation}
  \sups |\ghat_j(t) - \gamma_j(t)|
  = \Op(n^{-1/2}) + \Op(\tau_m^{1/2}) + \mathcal{O}(\eta_1^{1/2}),
  \qquad j = 1,\ldots,p,
  \label{eq:gamma-rate}
\end{equation}
and the same rate for the global alignment estimator of $\lambda$. Here
$\tau_m = m^{-(1-\delta)/3}$ for an arbitrarily small $\delta > 0$, $m$ is the number of
observations per curve, and $\eta_1$ is the penalty parameter of the pairwise warping step.

The global alignment estimator uses one of the $p$ observed curves per subject and discards
the rest. A natural alternative is to estimate $\lambda$ from the $p$ estimated tempos
$\ghat_1,\ldots,\ghat_p$, each of which already pools all $n$ subjects. Since the tempos
are warped versions of one curve, $\lambda$ is their mean in a space where curves are compared
up to time warping, that is, a Fr\'echet (or Karcher) mean. The standardization
\eqref{eq:std-psi} then selects which warped version of that mean is $\lambda$. Consistency of
an estimator of this type was not established in Carroll and M\"{u}ller (2023). In this note
we show that a Fr\'echet mean of the estimated tempos under the uniform metric, centered
according to \eqref{eq:std-psi}, attains the rate in \eqref{eq:gamma-rate}.

The organization of this note is as follows. Section \ref{sec:est} defines the estimator.
Section \ref{sec:results} states the two additional assumptions and the results, with a
plain-language overview of the argument. Section \ref{sec:proofs} contains the proofs.

\section{The Fr\'echet Mean Estimator}
\label{sec:est}

Let $\W$ denote the set of all strictly increasing continuous bijections of $\T$ onto
itself. This set contains the deformation space of Carroll and M\"{u}ller (2023) and is
closed under composition, inversion, and convex combinations. For continuous functions
$f, g$ on $\T$, define the warping-invariant distance
\begin{equation}
  d(f, g) = \inf_{\varphi \in \W} \norm{f \circ \varphi - g}.
  \label{eq:d}
\end{equation}
Since $\norm{f\circ\varphi - g\circ\varphi} = \norm{f - g}$ for every $\varphi \in \W$, the
function $d$ is symmetric, satisfies the triangle inequality, and satisfies
$d(f,g) \le \norm{f-g}$ and $d(f\circ\varphi, g) = d(f,g)$. It is a pseudometric on $C(\T)$
under which each tempo is at distance zero from the latent curve: $d(\gamma_j, \lambda) = 0$,
with the infimum attained at $\varphi = \Psi_j^{-1}$.

We assume the estimated tempos $\ghat_j$ in \eqref{eq:gammahat} are continuous. The sample Fr\'echet functional is
\[
  \hat{V}(\mu) = \frac{1}{p}\sum_{j=1}^p d^2(\ghat_j, \mu), \qquad \mu \in C(\T).
\]
A minimizer of $\hat{V}$ is determined at most up to time warping, and we do not assume that
the infimum is attained. We therefore allow a tolerance $\rho \ge 0$ in the two optimization
steps. The Fr\'echet mean estimator (FME) of $\lambda$ is obtained as follows.

\begin{enumerate}
  \item[Step 1.] (Template) Choose $\muhat \in C(\T)$ such that
    $\hat{V}(\muhat) \le \inf_{\mu \in C(\T)} \hat{V}(\mu) + \rho^2$.
  \item[Step 2.] (Alignment) For $j = 1,\ldots,p$, choose $\phihat_j \in \W$ such that
    $\norm{\ghat_j\circ\phihat_j - \muhat} \le d(\ghat_j, \muhat) + \rho$.
  \item[Step 3.] (Centering) Set $\bar{\varphi} = p^{-1}\sum_{j=1}^p \phihat_j$ and
    \begin{equation}
      \lhat_F = \muhat \circ \bar{\varphi}^{-1}.
      \label{eq:fme}
    \end{equation}
\end{enumerate}

For every $\rho > 0$, Steps 1 and 2 can be carried out by the definition of the infimum.
The centering step is the sample version of \eqref{eq:std-psi}: the warping functions
$\phihat_j \circ \bar{\varphi}^{-1}$ that carry the estimated tempos to $\lhat_F$ average
to the identity, just as the functions $\Psi_j^{-1}$ that carry the tempos to $\lambda$ do.

\section{Theoretical Results}
\label{sec:results}

In addition to (L1), (L2), and (S0--S2), we require the following.

\begin{itemize}
  \item[(L3)] The derivative $\lambda'$ has at most finitely many zeros in $\T$, and
    $\lambda''(z) \neq 0$ at each zero $z$.
  \item[(D)] For $j = 1,\ldots,p$, the inverse deformation $\Psi_j^{-1}$ is Lipschitz
    continuous on $\T$ with constant $L$.
\end{itemize}

Assumption (L3) strengthens the requirement in (L1) that $\lambda$ has no flat stretches: the
latent curve has finitely many peaks and valleys, and each has nonzero curvature. Assumption
(D) holds whenever the derivatives $\Psi_j'$ are bounded away from zero.

The argument has three parts. First, since $\lambda$ has Fr\'echet
variance at most $\max_j\norm{\ghat_j-\gamma_j}^2$, the triangle inequality for $d$ forces any
approximate minimizer $\muhat$ to be close to a warped version of $\lambda$. No curvature
condition on the Fr\'echet functional is needed, because the population Fr\'echet variance is
zero. Second, we need to control the warping functions. If $\lambda\circ g$ is uniformly
close to $\lambda$, then $g$ is close to the identity, but only at a square-root rate near the
peaks and valleys of $\lambda$, where the curve is flat to first order. Third, it emerges
that this loss does not reach the estimator: the warping error is large only where
$\lambda'$ is small, so its effect on the values of $\lambda$ is still of first order. This
is the content of Lemma \ref{lem:stab}, which bounds the product $|\lambda'(t)|\,|g(t)-t|$
rather than $|g(t)-t|$ alone.

Throughout, let $Z$ be the set of zeros of $\lambda'$, $K$ its cardinality,
$N \le K+1$ the number of intervals into which $Z \cup \{0,1\}$ partitions $\T$, and
$B = \norm{\lambda''}$. On each of the $N$ intervals $\lambda$ is strictly monotone.

\begin{lemma}[Stability of warping]
  \label{lem:stab}
  Suppose $\lambda \in C^2(\T)$ satisfies (L3). There exists a constant $c_0 > 0$ depending
  only on $\lambda$ such that
  \begin{equation}
    \int_s^t |\lambda'(u)|\,du \ge c_0\,(t-s)^2 \qquad \text{for all } 0 \le s \le t \le 1.
    \label{eq:quad}
  \end{equation}
  Moreover, if $g \in \W$ and $\delta \ge 0$ satisfy $\norm{\lambda\circ g - \lambda} \le \delta$,
  then for all $t \in \T$,
  \begin{equation}
    (g(t)-t)^2 \le C_\lambda\,\delta
    \qquad\text{and}\qquad
    |\lambda'(t)|\,|g(t)-t| \le C_\lambda\,\delta,
    \label{eq:stab}
  \end{equation}
  where $C_\lambda = 4N\max\{1, (1+B)/c_0\}$.
\end{lemma}

In particular, $\lambda \circ g = \lambda$ with $g \in \W$ implies $g = \id$.

\begin{theorem}[Deterministic bound]
  \label{thm:det}
  Suppose $\lambda \in C^2(\T)$ satisfies (L3), the deformations satisfy
  \eqref{eq:std-psi} and (D), and $\ghat_1,\ldots,\ghat_p \in C(\T)$. Let
  $\varepsilon = \max_{1\le j\le p}\norm{\ghat_j - \gamma_j}$, and let $\lhat_F$ be obtained
  from Steps 1--3 with tolerance $\rho \ge 0$. Then
  \begin{equation}
    \norm{\lhat_F - \lambda} \le C\,(\varepsilon + \rho),
    \qquad
    C = (1+\sqrt{p})\,\bigl\{1 + 2\,C_\lambda L\,(1 + BL/2)\bigr\}.
    \label{eq:det}
  \end{equation}
\end{theorem}

The bound holds for every choice of $\muhat$ and $\phihat_j$ that meets the tolerance, so
the estimator does not need to be unique. The case $\varepsilon = \rho = 0$ gives
identification at the population level.

\begin{corollary}[Population identification]
  \label{cor:pop}
  Under the assumptions of Theorem \ref{thm:det}, suppose $\ghat_j = \gamma_j$ for all $j$.
  Then $\mu = \lambda$ and $\varphi_j = \Psi_j^{-1}$ satisfy Steps 1 and 2 with $\rho = 0$, and
  every output of Steps 1--3 with $\rho = 0$ is equal to $\lambda$.
\end{corollary}

\begin{theorem}[Rate of convergence]
  \label{thm:rate}
  Suppose assumptions (L1), (L2), and (S0--S2) of Carroll and M\"{u}ller (2023) hold together
  with (L3) and (D), and let $\tau_m = m^{-(1-\delta)/3}$ for an arbitrarily small $\delta > 0$.
  If the FME is computed from the estimated tempos $\ghat_j$ with tolerance $\rho_n$, then
  \[
    \sups |\lhat_F(t) - \lambda(t)|
    = \Op(n^{-1/2}) + \Op(\tau_m^{1/2}) + \mathcal{O}(\eta_1^{1/2}) + \mathcal{O}(\rho_n).
  \]
\end{theorem}

\begin{corollary}
  \label{cor:param}
  Under the assumptions of Theorem \ref{thm:rate}, if additionally
  $\eta_1 = \mathcal{O}(n^{-1})$, $\rho_n = \mathcal{O}(n^{-1/2})$, and $m$ is bounded below
  by a multiple of $n^{\Delta(1-\delta)^{-1}}$ for some $\Delta \ge 3$, then
  \[
    \sups |\lhat_F(t) - \lambda(t)| = \Op(n^{-1/2}).
  \]
\end{corollary}

The FME therefore attains the same rate as the global alignment estimator, including the
parametric rate on dense enough designs, while using all $p$ components of every subject.

\section{Proofs}
\label{sec:proofs}

\begin{proof}[Proof of Lemma \ref{lem:stab}]
We first show \eqref{eq:quad}. If $Z$ is empty, then $c_1 := \min_{u\in\T}|\lambda'(u)| > 0$ and
$\int_s^t|\lambda'| \ge c_1(t-s) \ge c_1(t-s)^2$. Otherwise, write $\mathrm{dist}(u,Z)$ for the
distance from $u$ to the nearest zero of $\lambda'$ and consider the ratio
$r(u) = |\lambda'(u)|/\mathrm{dist}(u,Z)$ on $\T\setminus Z$. The ratio is continuous and
positive, and since $\lambda'(u) = \lambda''(z)(u-z) + o(|u-z|)$ as $u \to z$, it converges to
$|\lambda''(z)| > 0$ at each $z \in Z$. Hence $r$ extends to a positive continuous function
on the compact set $\T$, and $c_1 := \inf_u r(u) > 0$, that is,
$|\lambda'(u)| \ge c_1\,\mathrm{dist}(u,Z)$ on $\T$. Now fix $s < t$ and let $\ell = t-s$. The
set of $u\in[s,t]$ with $\mathrm{dist}(u,Z) < \ell/(4K)$ has Lebesgue measure at most
$2K\cdot\ell/(4K) = \ell/2$, so on a subset of $[s,t]$ of measure at least $\ell/2$ we have
$|\lambda'(u)| \ge c_1\ell/(4K)$. Therefore $\int_s^t|\lambda'| \ge c_1\ell^2/(8K)$, and
\eqref{eq:quad} holds in both cases with $c_0 = c_1/(8\max\{K,1\})$.

Next, let $V(t) = \int_0^t|\lambda'(u)|\,du$, the total variation of $\lambda$ on $[0,t]$, and
let $h = \lambda\circ g$. We claim that
\begin{equation}
  |V(g(t)) - V(t)| \le 2N\delta \qquad \text{for all } t\in\T.
  \label{eq:tv}
\end{equation}
Let $0 = a_0 < a_1 < \cdots < a_N = 1$ be the points of $Z\cup\{0,1\}$, so that $\lambda$ is
monotone on each $[a_{i-1},a_i]$, and let $q_i = g^{-1}(a_i)$, so that $h$ is monotone on
each $[q_{i-1},q_i]$. Since $g$ maps $[0,t]$ increasingly onto $[0,g(t)]$, the total
variation of $h$ on $[0,t]$ is $V(g(t))$. For a function that is monotone on each interval of
a partition, the total variation on $[0,t]$ is the sum of the absolute increments over the
partition truncated at $t$, and for any function the total variation is at least that sum.
Applying this to $h$ with the partition $\{q_i\}$, and using $\norm{h-\lambda}\le\delta$ on
each of the at most $N$ nonzero increments,
\[
  V(g(t)) = \sum_{i=1}^N |h(q_i\wedge t) - h(q_{i-1}\wedge t)|
  \le \sum_{i=1}^N |\lambda(q_i\wedge t) - \lambda(q_{i-1}\wedge t)| + 2N\delta
  \le V(t) + 2N\delta.
\]
Exchanging the roles of $\lambda$ and $h$ and using the partition $\{a_i\}$ gives
$V(t) \le V(g(t)) + 2N\delta$, which proves \eqref{eq:tv}.

Fix $t$ and let $\ell = |g(t)-t|$. By \eqref{eq:tv} and \eqref{eq:quad},
\[
  c_0\,\ell^2 \le \Bigl|\int_t^{g(t)}|\lambda'(u)|\,du\Bigr| = |V(g(t)) - V(t)| \le 2N\delta,
\]
so $\ell^2 \le 2N\delta/c_0 \le C_\lambda\delta$, which is the first bound in \eqref{eq:stab}.
For the second, if $2B\ell \le |\lambda'(t)|$, then $|\lambda'(u)| \ge |\lambda'(t)| - B\ell
\ge |\lambda'(t)|/2$ for all $u$ between $t$ and $g(t)$, and \eqref{eq:tv} gives
$|\lambda'(t)|\,\ell/2 \le 2N\delta$, that is, $|\lambda'(t)|\,\ell \le 4N\delta$. If instead
$|\lambda'(t)| < 2B\ell$, then $|\lambda'(t)|\,\ell < 2B\ell^2 \le 4BN\delta/c_0$. In both cases
$|\lambda'(t)|\,\ell \le C_\lambda\delta$.
\end{proof}

\begin{proof}[Proof of Theorem \ref{thm:det}]
Let $\delta_0 = (1+\sqrt{p})(\varepsilon+\rho)$.

\emph{Step (a): the template is close to the tempos.} Since $d(\gamma_j,\lambda) = 0$ and
$d(\ghat_j,\gamma_j) \le \norm{\ghat_j-\gamma_j} \le \varepsilon$, the triangle inequality gives
$d(\ghat_j,\lambda)\le\varepsilon$ for all $j$, hence $\hat{V}(\lambda)\le\varepsilon^2$ and,
by Step 1 of the estimator, $\hat{V}(\muhat) \le \varepsilon^2+\rho^2$. Each summand of
$\hat{V}(\muhat)$ is then at most $p(\varepsilon^2+\rho^2)$, so
$d(\ghat_j,\muhat) \le \sqrt{p}\,(\varepsilon+\rho)$. By Step 2 of the estimator and the
invariance of the uniform norm under warping,
\begin{equation}
  \norm{\gamma_j\circ\phihat_j - \muhat}
  \le \norm{\gamma_j - \ghat_j} + \norm{\ghat_j\circ\phihat_j - \muhat}
  \le \varepsilon + \sqrt{p}\,(\varepsilon+\rho) + \rho = \delta_0,
  \qquad j = 1,\ldots,p.
  \label{eq:a}
\end{equation}

\emph{Step (b): the alignment warps are stable.} Let $\hat{\rho} = \Psi_1\circ\phihat_1 \in \W$
and $g_j = \Psi_j\circ\phihat_j\circ\hat{\rho}^{-1} \in \W$, so that $g_1 = \id$. Since
$\gamma_j\circ\phihat_j = \lambda\circ\Psi_j\circ\phihat_j$, applying \eqref{eq:a} for $j$ and
for $1$ gives $\norm{\lambda\circ\Psi_j\circ\phihat_j - \lambda\circ\hat{\rho}} \le 2\delta_0$,
and composing with $\hat{\rho}^{-1}$,
\[
  \norm{\lambda\circ g_j - \lambda} \le 2\delta_0, \qquad j = 1,\ldots,p.
\]
Lemma \ref{lem:stab} then yields, for all $t\in\T$ and all $j$,
\begin{equation}
  (g_j(t)-t)^2 \le 2C_\lambda\delta_0
  \qquad\text{and}\qquad
  |\lambda'(t)|\,|g_j(t)-t| \le 2C_\lambda\delta_0.
  \label{eq:b}
\end{equation}

\emph{Step (c): centering.} By construction $\phihat_j = \Psi_j^{-1}\circ g_j\circ\hat{\rho}$, so
$\bar{\varphi} = \kappa\circ\hat{\rho}$ with
$\kappa = p^{-1}\sum_{j=1}^p \Psi_j^{-1}\circ g_j \in \W$. The standardization
\eqref{eq:std-psi} and (D) give
\begin{equation}
  |\kappa(t) - t|
  = \Bigl|\frac{1}{p}\sum_{j=1}^p \bigl\{\Psi_j^{-1}(g_j(t)) - \Psi_j^{-1}(t)\bigr\}\Bigr|
  \le L\max_{1\le j\le p}|g_j(t)-t|.
  \label{eq:kappa}
\end{equation}
Since $\bar{\varphi}^{-1} = \hat{\rho}^{-1}\circ\kappa^{-1}$, the definition \eqref{eq:fme} and
\eqref{eq:a} with $j=1$ give
\begin{align}
  \norm{\lhat_F - \lambda}
  &\le \norm{\muhat\circ\hat{\rho}^{-1}\circ\kappa^{-1} - \lambda\circ\kappa^{-1}}
     + \norm{\lambda\circ\kappa^{-1} - \lambda} \nonumber\\
  &= \norm{\muhat - \lambda\circ\hat{\rho}} + \norm{\lambda - \lambda\circ\kappa}
  \le \delta_0 + \norm{\lambda\circ\kappa - \lambda}.
  \label{eq:c}
\end{align}

\emph{Step (d): first-order effect of the centering error.} A second-order Taylor expansion
of $\lambda$ at $t$, together with \eqref{eq:kappa} and \eqref{eq:b}, gives
\begin{align*}
  |\lambda(\kappa(t)) - \lambda(t)|
  &\le |\lambda'(t)|\,|\kappa(t)-t| + \tfrac{1}{2}B\,(\kappa(t)-t)^2\\
  &\le L\max_j |\lambda'(t)|\,|g_j(t)-t| + \tfrac{1}{2}BL^2\max_j (g_j(t)-t)^2\\
  &\le 2C_\lambda L\,(1 + BL/2)\,\delta_0
\end{align*}
for all $t \in \T$. Combining this with \eqref{eq:c} proves \eqref{eq:det}.
\end{proof}

\begin{proof}[Proof of Corollary \ref{cor:pop}]
With $\ghat_j = \gamma_j$ we have $\hat{V}(\lambda) = 0$, so $\mu = \lambda$ satisfies Step 1
with $\rho = 0$, and $\gamma_j\circ\Psi_j^{-1} = \lambda$ shows that $\varphi_j = \Psi_j^{-1}$
satisfies Step 2 with $\rho = 0$. Theorem \ref{thm:det} with $\varepsilon = \rho = 0$ gives
$\norm{\lhat_F - \lambda} = 0$ for every output of Steps 1--3.
\end{proof}

\begin{proof}[Proof of Theorem \ref{thm:rate}]
The constant $C$ in \eqref{eq:det} depends only on $\lambda$, $L$, and $p$, and
$\varepsilon = \max_j\norm{\ghat_j-\gamma_j}$ has the rate in \eqref{eq:gamma-rate} by
Theorem 1b of Carroll and M\"{u}ller (2023), since $p$ is fixed. The result follows from
Theorem \ref{thm:det}.
\end{proof}

\begin{proof}[Proof of Corollary \ref{cor:param}]
If $m \ge c\,n^{\Delta(1-\delta)^{-1}}$ with $\Delta \ge 3$, then
$\tau_m^{1/2} = m^{-(1-\delta)/6} = \mathcal{O}(n^{-\Delta/6}) = \mathcal{O}(n^{-1/2})$. The
remaining terms are $\mathcal{O}(n^{-1/2})$ by assumption.
\end{proof}

\bigskip
\noindent\textbf{References}\\[4pt]
Carroll, C. \& M\"{u}ller, H.-G. (2023). Latent deformation models for
multivariate functional data and time-warping separability.
\textit{Biometrics}, 79(4), 3345--3358.\\[4pt]
Tang, R. \& M\"{u}ller, H.-G. (2008). Pairwise curve synchronization for functional data.
\textit{Biometrika}, 95(4), 875--889.

\end{document}
